%%%%%%%%%%%%%%%%%%%%%%%%%%%%%%%%%%%%%%%%%%%%%%%%%%%%%%%%%%%%%%%%%%%%%%%%%%%%%%%%
% Preámbulo                                                                    %
%%%%%%%%%%%%%%%%%%%%%%%%%%%%%%%%%%%%%%%%%%%%%%%%%%%%%%%%%%%%%%%%%%%%%%%%%%%%%%%%

\documentclass[11pt,a4paper,titlepage,oneside]{report}

%%% RELACIÓN DE VARIABLES A PERSONALIZAR %%%
%%\def\lingua{gal}
\def\lingua{esp} % descomenta esta liña se redactarás a memoria en español
%\def\lingua{eng} % descomenta esta liña se redactarás a memoria en inglés
\def\nome{Pablo Soage Rodas}                             % substitúe aquí o teu nome
\def\nomedirectorA{Enrique Fernández Blanco}              % substitúe aquí o nome de quen dirixe
\def\nomedirectorB{Daniel Rivero Cebrián}             % duplica esta liña máis veces se o precisas, cambiando
                                                     % a letra final (A, B, C, D...): úsanse na portada.tex
\def\titulo{Framework de Captura e Sincronización de Datos para a Creación de Datasets Orientados a Interfaces Cerebro-Computador} % substitúe aquí o título do teu TFG
%\def\titulacion{gced}                               % descomenta esta liña e comenta a seguinte se es estudante do GCED
\def\titulacion{gei}
\def\mencion{COMPUTACIÓN}                           % descomenta a mención que che corresponda se es estudante do GEI
%\def\mencion{ENXEÑARÍA DO SOFTWARE}
%\def\mencion{ENXEÑARÍA DE COMPUTADORES}
%\def\mencion{SISTEMAS DE INFORMACIÓN}
%\def\mencion{TECNOLOXÍAS DA INFORMACIÓN}

\def\renomearcadros{si} % descomenta esta liña se redactas a memoria en español e prefires que
                         % os "cuadros" e o "índice de cuadros" se renomeen
                         % a "tablas" e "índice de tablas" respectivamente

\usepackage{estilo_tfg}

% Lista de paquetes potencialmente interesantes (uso baixo demanda)

% \usepackage{alltt}       % proporciona o entorno alltt, semellante a verbatim pero que respecta comandos
% \usepackage{enumitem}    % permite personalizar os entornos de lista
% \usepackage{eurofont}    % proporciona o comando \euro
\usepackage{float}       % permite máis opcións para controlar obxectos flotantes (táboas, figuras)
% \usepackage{hhline}      % permite personalizar as liñas horizontais en arrays e táboas
  \usepackage{longtable}   % permite construir táboas que ocupan máis dunha páxina
% \usepackage{lscape}      % permite colocar partes do documento en orientación apaisada
% \usepackage{moreverb}    % permite personalizar o entorno verbatim
% \usepackage{multirow}    % permite crear celdas que ocupan varias filas da mesma táboa
% \usepackage{pdfpages}    % permite insertar ficheiros en PDF no documento
% \usepackage{rotating}    % permite diferentes tipos de rotacións para figuras e táboas
\usepackage{subcaption}  % permite a inclusión de varias subfiguras nunha figura
% \usepackage{tabu}        % permite táboas flexibles
% \usepackage{tabularx}    % permite táboas con columnas de anchura determinada
\usepackage{listings}    % permite incluir código fuente con resaltado de sintaxis

%%%%%%%%%%%%%%%%%%%%%%%%%%%%%%%%%%%%%%%%%%%%%%%%%%%%%%%%%%%%%%%%%%%%%%%%%%%%%%%%
% Corpo                                                                        %
%%%%%%%%%%%%%%%%%%%%%%%%%%%%%%%%%%%%%%%%%%%%%%%%%%%%%%%%%%%%%%%%%%%%%%%%%%%%%%%%

\begin{document}

 %%%%%%%%%%%%%%%%%%%%%%%%%%%%%%%%%%%%%%%%
 % Preliminares do documento            %
 %%%%%%%%%%%%%%%%%%%%%%%%%%%%%%%%%%%%%%%%

 \begin{titlepage}

  \hspace*{128pt}
  \textcolor{udcpink}{{\fontencoding{T1}\fontfamily{phv}\selectfont Facultade de Informática}}\\[-32pt]

  \begin{center}
    \includegraphics[scale=0.3]{imaxes/udc}\\[25pt]

    {\large TRABALLO FIN DE GRAO \\
            \nometitulacion \\
            \nomemencion } \\[10pt]

    \carimbo \\[25pt]

    \begin{huge}
      \begin{spacing}{1.3}
        \bfseries \titulo
      \end{spacing}
    \end{huge}
  \end{center}

  \vfill

  \begin{flushright}
    {\large
    \begin{tabular}{ll}
      {\bf Estudante:} & \nome \\
      {\bf Dirección:} & \nomedirectorA \\
                      & \nomedirectorB \\ % duplica esta liña máis veces se o precisas, cambiando
                                           % a letra final (A, B, C, D...); define eses nomes no memoria_tfg.tex
    \end{tabular}}
  \end{flushright}
  \rightline{A Coruña, \datasimple.}
\end{titlepage}

 \dedicatoria{A mis padres, por su apoyo incondicional durante todos mis años de estudio.} % escribe neste comando o teu texto de dedicatoria
 \paxinaenbranco
 \begin{agradecementos}
 Quiero expresar mi más sincero agradecimiento a mis padres por apoyarme en lo que fuese necesario durante este período de estudio, llevarme y traerme del bus y preocuparse siempre por lo que nos pudiese hacer falta.
 A mis abuelos, por estar cuidando constantemente de nosotros estos años.
 Reconocimiento especial a mi abuelo Manuel Rodas y a alguien que para mi fue como un abuelo, Crisanto Domínguez, que ahora ya no están a mi lado, por haberme inculcado desde pequeño la curiosidad por todas las cosas y las ansias por aprender.
 A mi pareja, Jennifer Freire, por haber estado ahí siempre, a mi lado, apoyandome incluso cuando la situación era muy dura.
 A mi hermana, que siempre ha estado orgullosa de mi.
 A mis amigos, que han estado conmigo a lo largo de este período de estudio, trabajando a mi lado y haciéndome reír en momentos no favorables.
 \\\\También agradezco a la Facultad de Informática de la Universidade da Coruña por proporcionarme acceso al dispositivo g.tec Unicorn Hybrid Black, así como a mis tutores de TFG y el director del CITIC por confiarme dicho dispositivo.
 \end{agradecementos}
 %%%%%%%%%%%%%%%%%%%%%%%%%%%%%%%%%%%%%%%%%%%%%%%%%%%%%%%%%%%%%%%%%%%%%%%%%%%%%%%%

\pagestyle{empty}
\begin{abstract}
  Las interfaces cerebro-computadora (BCI) representan un paradigma prometedor para la interacción humano-máquina, pero su desarrollo enfrenta limitaciones significativas debido a la falta de herramientas accesibles que faciliten la captura sincronizada de señales electroencefalográficas (EEG) junto con eventos de entrada convencionales del usuario. Esta sincronización precisa es fundamental para establecer correlaciones entre actividad cerebral y acciones del usuario, requisito esencial para el entrenamiento de modelos predictivos en aplicaciones BCI.

  Este trabajo presenta el desarrollo de un \textit{framework} modular de código abierto que aborda esta problemática integrando la captura sincronizada de señales EEG y eventos de entrada de usuario (teclado, ratón, controladores) con precisión temporal de milisegundos. El sistema implementa una arquitectura basada en principios SOLID con patrones dirigidos por eventos, proporcionando múltiples interfaces de usuario (GUI PyQt5, CLI y API REST con WebSocket) que facilitan tanto el uso interactivo como la integración programática.

  La solución soporta dispositivos comerciales de EEG mediante BrainFlow, incluyendo el g.tec Unicorn Hybrid Black, e incorpora procesamiento de señales en tiempo real con algoritmos avanzados como análisis de componentes independientes (ICA), transformadas \textit{wavelet} y filtros digitales. El \textit{framework} incluye visualización de señales biomédicas, monitoreo de calidad de señal, y exportación a formatos estándar como BIDS, EDF+, HDF5 y CSV.

  La validación del sistema con datos sintéticos demuestra sincronización precisa entre señales EEG y eventos de entrada, cumpliendo los requisitos para la generación de \textit{datasets} de calidad orientados al desarrollo de interfaces cerebro-computadora. El \textit{framework} desarrollado constituye una herramienta completa que integra capacidades tradicionalmente fragmentadas en el ecosistema BCI, facilitando la investigación y desarrollo en este campo emergente.

  \vspace*{25pt}
  \begin{segundoresumo}
    Brain-computer interfaces (BCI) represent a promising paradigm for human-machine interaction, but their development faces significant limitations due to the lack of accessible tools that facilitate synchronized capture of electroencephalographic (EEG) signals along with conventional user input events. This precise synchronization is fundamental for establishing correlations between brain activity and user actions, an essential requirement for training predictive models in BCI applications.

    This work presents the development of an open-source modular \textit{framework} that addresses this problem by integrating synchronized capture of EEG signals and user input events (keyboard, mouse, controllers) with millisecond temporal precision. The system implements an architecture based on SOLID principles with event-driven patterns, providing multiple user interfaces (PyQt5 GUI, CLI, and REST API with WebSocket) that facilitate both interactive use and programmatic integration.

    The solution supports commercial EEG devices through BrainFlow, including the g.tec Unicorn Hybrid Black, and incorporates real-time signal processing with advanced algorithms such as Independent Component Analysis (ICA), \textit{wavelet} transforms, and digital filters. The \textit{framework} includes biomedical signal visualization, signal quality monitoring, and export to standard formats such as BIDS, EDF+, HDF5, and CSV.

    System validation with synthetic data demonstrates precise synchronization between EEG signals and input events, meeting requirements for generating quality \textit{datasets} oriented towards brain-computer interface development. The developed \textit{framework} constitutes a comprehensive tool that integrates traditionally fragmented capabilities in the BCI ecosystem, facilitating research and development in this emerging field.
  \end{segundoresumo}
\vspace*{25pt}
\begin{multicols}{2}
\begin{description}
\item [\palabraschaveprincipal:] \mbox{} \\[-20pt]
  \begin{itemize}
    \item Interfaces cerebro-computadora (BCI)
    \item Electroencefalografía (EEG)
    \item Sincronización de datos
    \item Framework modular
    \item Procesamiento de señales
    \item Captura de eventos
    \item Datasets biomédicos
    \item BrainFlow
    \item Arquitectura dirigida por eventos
    \item API REST
    \item PyQt5
    \item Python
  \end{itemize}
\end{description}
\begin{description}
\item [\palabraschavesecundaria:] \mbox{} \\[-20pt]
  \begin{itemize}
    \item Brain-computer interfaces (BCI)
    \item Electroencephalography (EEG)
    \item Data synchronization
    \item Modular framework
    \item Signal processing
    \item Event capture
    \item Biomedical datasets
    \item BrainFlow
    \item Event-driven architecture
    \item REST API
    \item PyQt5
    \item Python
  \end{itemize}
\end{description}
\end{multicols}

\end{abstract}
\pagestyle{fancy}

%%%%%%%%%%%%%%%%%%%%%%%%%%%%%%%%%%%%%%%%%%%%%%%%%%%%%%%%%%%%%%%%%%%%%%%%%%%%%%%%

 \pagenumbering{roman}
 \setcounter{page}{1}
 \bstctlcite{IEEEexample:BSTcontrol}

 \tableofcontents
 \listoffigures
 \listoftables
 \clearpage
 
 \pagenumbering{arabic}
 \setcounter{page}{1}

 %%%%%%%%%%%%%%%%%%%%%%%%%%%%%%%%%%%%%%%%
 % Capítulos                            %
 %%%%%%%%%%%%%%%%%%%%%%%%%%%%%%%%%%%%%%%%

 \chapter{Introducción}
\label{chap:introduccion}

\lettrine{L}{a} evolución de la interacción humano-computadora ha experimentado transformaciones revolucionarias en las últimas décadas. Desde las primeras interfaces basadas en comandos de texto hasta los modernos sistemas táctiles y de reconocimiento de voz, cada avance ha ampliado las posibilidades de comunicación entre personas y máquinas~\cite{eeg_bci_review,wolpaw2002brain}. Sin embargo, todas estas modalidades comparten una característica común: requieren algún tipo de actividad física consciente por parte del usuario, ya sea escribir, hacer clic, tocar una pantalla o hablar.

Las \acrfull{bci} representan un paradigma fundamentalmente diferente que podría revolucionar esta interacción~\cite{eeg_bci_review,wolpaw2002brain}. Estos sistemas permiten establecer canales de comunicación directa entre el cerebro humano y dispositivos electrónicos, interpretando las señales neurológicas sin necesidad de actividad muscular consciente~\cite{dornhege2007toward}. Esta capacidad abre posibilidades extraordinarias en campos tan diversos como la medicina rehabilitadora, donde pueden restaurar capacidades de comunicación a personas con lesiones medulares severas~\cite{wolpaw2002brain}, la accesibilidad tecnológica, proporcionando nuevas formas de interacción para personas con discapacidades motoras, y el desarrollo de paradigmas de interacción completamente nuevos que podrían transformar la forma en que los humanos utilizamos la tecnología.

El fundamento de los sistemas \acrshort{bci} reside en la capacidad de capturar y interpretar \gls{biosenial} generadas por la actividad neuronal. Los sistemas \acrshort{bci} modernos se clasifican principalmente según el método de adquisición de señales neurológicas, incluyendo la \acrfull{eeg}, la \acrfull{meg}, la \acrfull{fnirs} y técnicas invasivas como la \acrfull{ecog}. Entre estas modalidades, la \acrshort{eeg} se ha establecido como la técnica predominante para aplicaciones \acrshort{bci} no invasivas debido a su excelente resolución temporal (del orden de milisegundos), su naturaleza completamente no invasiva y la disponibilidad creciente de dispositivos comerciales accesibles~\cite{eeg_bci_review,wolpaw2002brain,dornhege2007toward}. A continuación se analizará en detalle el panorama actual de estas tecnologías y las limitaciones que motivan el desarrollo de nuevas soluciones integradas.

Uno de los obstáculos más importantes en el desarrollo de sistemas \acrshort{bci} efectivos es la necesidad de \gls{dataset} de alta calidad que contengan señales \acrshort{eeg} sincronizadas temporalmente con las intenciones o acciones del usuario. Esta \gls{sincronizacion} es crítica porque permite establecer correlaciones estadísticamente significativas entre patrones específicos de actividad cerebral y las intenciones del usuario~\cite{eeg_desynchronization}. El fundamento neurofisiológico de estas correlaciones reside en fenómenos como la \gls{desincronizacion}, donde la actividad neuronal de fondo se modifica cuando el cerebro se prepara para una acción específica. Por ejemplo, cuando una persona imagina mover su mano derecha, se produce una disminución en la potencia de las frecuencias mu (8-13Hz) y beta (13-30Hz) sobre las áreas motoras contralaterales del cerebro~\cite{eeg_desynchronization}. Capturar precisamente cuándo ocurre esta imaginación de movimiento en relación con el patrón \acrshort{eeg} observado es fundamental para entrenar algoritmos de \acrfull{ml} que puedan reconocer y clasificar estas intenciones en tiempo real. Sin esta sincronización temporal precisa, los modelos predictivos no pueden establecer las asociaciones necesarias entre estados cerebrales y intenciones del usuario, limitando significativamente la efectividad de los sistemas \acrshort{bci}.

Sin embargo, el ecosistema actual de herramientas para el desarrollo de sistemas \acrshort{bci} presenta una fragmentación significativa que dificulta la creación de estos \gls{dataset} de calidad. Las herramientas existentes tienden a abordar aspectos específicos del problema de forma aislada, sin proporcionar soluciones integradas que cubran todo el flujo de trabajo necesario desde la captura hasta el análisis de datos.

\section{Definición del problema}
\label{sec:problema}

El problema central que motiva este trabajo es la ausencia de una plataforma integrada que facilite la generación de \gls{dataset} de calidad para el desarrollo de sistemas \acrshort{bci}. Esta carencia se manifiesta en varios aspectos críticos que limitan tanto la investigación académica como el desarrollo de aplicaciones prácticas en el campo.

En primer lugar, no existe una solución que proporcione captura sincronizada de señales \acrshort{eeg} junto con eventos de entrada convencionales del usuario. Esta \gls{sincronizacion} temporal es fundamental para establecer las correlaciones necesarias entre actividad cerebral e intenciones del usuario~\cite{eeg_desynchronization,Pontifex2023}. Las herramientas actuales como \textit{Lab Streaming Layer} requieren programación personalizada sustancial y conocimientos técnicos avanzados, mientras que bibliotecas como \textit{BrainFlow} se centran únicamente en la adquisición de datos \acrshort{eeg} sin proporcionar capacidades para la captura de eventos de usuario.

En segundo lugar, existe una carencia de interfaces de usuario accesibles que permitan a investigadores sin conocimientos de programación avanzados generar estos \gls{dataset}. La mayoría de herramientas existentes están diseñadas para desarrolladores expertos, creando barreras de entrada significativas para investigadores en neurociencia, psicología o medicina que podrían beneficiarse de estas tecnologías.

Adicionalmente, la falta de estandarización en formatos de exportación de datos complica el intercambio y reutilización de \gls{dataset} entre diferentes grupos de investigación. Cada herramienta tiende a utilizar formatos propietarios o específicos, dificultando la reproducibilidad de experimentos y la colaboración científica.

Por último, no existe una arquitectura extensible que permita integrar fácilmente nuevos tipos de dispositivos de entrada o algoritmos de procesamiento, limitando la evolución y adaptación de las herramientas a medida que avanzan las tecnologías \acrshort{bci}.

\section{Análisis \acrfull{dafo}}
\label{sec:analisis_dafo}

Para evaluar adecuadamente la viabilidad y posicionamiento estratégico de la solución propuesta, se presenta un análisis \acrshort{dafo} que examina las fortalezas internas del sistema, sus debilidades, las oportunidades del entorno y las amenazas externas que podrían afectar su desarrollo y adopción.

\subsection{Fortalezas del sistema propuesto}
\label{subsec:fortalezas}

\begin{itemize}
    \item \textbf{Integración multidispositivo}: El \textit{framework} proporciona por primera vez una solución unificada que integra múltiples fuentes de entrada (\acrshort{eeg}, teclado, ratón, controladores) en un único sistema sincronizado, eliminando la necesidad de combinar herramientas fragmentadas.

    \item \textbf{Arquitectura extensible y modular}: El diseño basado en principios \acrshort{solid} y patrones dirigidos por eventos facilita la incorporación de nuevos dispositivos, algoritmos de procesamiento y técnicas de análisis sin modificar el núcleo del sistema.

    \item \textbf{Múltiples interfaces de usuario}: La disponibilidad simultánea de interfaces gráficas, de línea de comandos y \acrshort{api} web permite que usuarios con diferentes niveles técnicos y contextos de uso puedan aprovechar las capacidades del sistema.

    \item \textbf{Soporte de formatos estándar}: La exportación nativa a formatos ampliamente utilizados en neurociencia (\acrshort{bids}, \acrshort{edf}+, \acrshort{hdf5}) aumenta significativamente la usabilidad para investigadores y facilita la interoperabilidad con herramientas existentes.

    \item \textbf{Procesamiento avanzado en tiempo real}: La implementación de algoritmos sofisticados como \acrshort{ica}, transformadas \textit{wavelet} y filtros digitales con optimizaciones para tiempo real proporciona capacidades de análisis que tradicionalmente requerían procesamiento \textit{offline}.
\end{itemize}

\subsection{Debilidades a considerar}
\label{subsec:debilidades}

\begin{itemize}
    \item \textbf{Dependencia de bibliotecas externas}: La dependencia crítica de bibliotecas como \textit{BrainFlow} para la adquisición de señales \acrshort{eeg} podría limitar las funcionalidades del sistema si estas bibliotecas presentan limitaciones o discontinúan su desarrollo.

    \item \textbf{Desafíos de sincronización temporal}: La sincronización precisa entre diferentes fuentes de datos presenta complejidades técnicas debido a latencias variables en diferentes sistemas operativos y configuraciones de hardware, lo que podría afectar la calidad de los \gls{dataset} generados.

    \item \textbf{Limitaciones de escalabilidad}: Aunque el sistema está optimizado para tiempo real, el procesamiento simultáneo de múltiples dispositivos \acrshort{eeg} de alta resolución podría requerir recursos computacionales significativos que limiten su aplicabilidad en sistemas con recursos limitados.

    \item \textbf{Curva de aprendizaje inicial}: A pesar de las interfaces simplificadas, el dominio \acrshort{bci} requiere conocimientos especializados en procesamiento de señales y neurociencia que podrían representar una barrera inicial para algunos usuarios.
\end{itemize}

\subsection{Oportunidades en el entorno}
\label{subsec:oportunidades}

\begin{itemize}
    \item \textbf{Crecimiento exponencial del campo \acrshort{bci}}: El interés creciente en aplicaciones \acrshort{bci} en medicina rehabilitadora, accesibilidad tecnológica y nuevas interfaces de usuario crea un mercado en expansión que demanda herramientas más accesibles y potentes.

    \item \textbf{Carencia de soluciones integradas}: La fragmentación actual del ecosistema de herramientas \acrshort{bci} representa una oportunidad significativa para una solución que integre captura sincronizada de señales \acrshort{eeg} y eventos de usuario.

    \item \textbf{Democratización de la investigación \acrshort{bci}}: La reducción de barreras técnicas puede facilitar el acceso a la investigación \acrshort{bci} para grupos con menos recursos técnicos, ampliando potencialmente la base de usuarios y contribuyendo al avance del campo.

    \item \textbf{Tendencia hacia código abierto}: La preferencia creciente por soluciones de código abierto en investigación científica posiciona favorablemente al \textit{framework} frente a alternativas comerciales cerradas.

    \item \textbf{Evolución de hardware accesible}: La disponibilidad creciente de dispositivos \acrshort{eeg} de bajo coste y alta calidad amplía el mercado potencial de usuarios que podrían beneficiarse del sistema.
\end{itemize}

\subsection{Amenazas externas}
\label{subsec:amenazas}

\begin{itemize}
    \item \textbf{Competencia de soluciones comerciales}: La aparición de nuevas soluciones comerciales similares, especialmente de empresas con mayores recursos, podría limitar la adopción del \textit{framework} de código abierto.

    \item \textbf{Evolución rápida del hardware \acrshort{eeg}}: Los avances acelerados en tecnología de dispositivos \acrshort{eeg} podrían requerir actualizaciones frecuentes y significativas del sistema para mantener la compatibilidad y relevancia.

    \item \textbf{Cambios en estándares de la industria}: La evolución de formatos de datos y protocolos de comunicación en el campo \acrshort{bci} podría requerir adaptaciones importantes del \textit{framework}.

    \item \textbf{Limitaciones regulatorias}: El desarrollo de regulaciones más estrictas para aplicaciones \acrshort{bci}, especialmente en contextos médicos, podría crear barreras adicionales para la adopción del sistema.

    \item \textbf{Fragmentación de la comunidad}: La falta de consenso en la comunidad \acrshort{bci} sobre estándares y mejores prácticas podría limitar la adopción uniforme de cualquier solución integrada.
\end{itemize}

\section{Objetivos}
\label{sec:objetivos}

Para abordar las limitaciones identificadas y aprovechar las oportunidades reveladas por el análisis \acrshort{dafo}, este trabajo define los siguientes objetivos que guiarán el desarrollo de la solución propuesta.

\subsection{Objetivo general}
\label{subsec:objetivo_general}

El objetivo principal de este \acrfull{tfg} es crear una plataforma software que facilite la captura, sincronización y procesamiento de señales \acrshort{eeg} junto con eventos de entrada de usuario, proporcionando las herramientas necesarias para la generación de \gls{dataset} de calidad que puedan ser utilizados en el desarrollo de \gls{bci}.

\subsection{Objetivos específicos}
\label{subsec:objetivos_especificos}

Del objetivo general se derivan los siguientes objetivos específicos que estructuran el desarrollo del proyecto:

\begin{enumerate}
    \item \textbf{Realizar un estudio exhaustivo del estado de la cuestión}: Analizar las herramientas existentes para el desarrollo de sistemas \acrshort{bci}, identificando sus fortalezas, limitaciones y oportunidades de mejora. Este análisis incluirá tanto sistemas existentes como bibliotecas de soporte en las que se puede basar el desarrollo.

    \item \textbf{Desarrollar un sistema robusto de captura de eventos de usuario}: Implementar un sistema que registre eventos de teclado, ratón y controladores con precisión temporal de milisegundos, incluyendo sistemas de marcado temporal para correlacionar estos eventos con señales cerebrales.

    \item \textbf{Implementar un sistema de adquisición y procesado de señales \acrshort{eeg}}: Desarrollar capacidades de captura de señales \acrshort{eeg} con soporte para dispositivos comerciales, incluyendo procesamiento de señales en tiempo real con algoritmos avanzados como filtrado digital, \acrfull{ica} y transformadas tiempo-frecuencia.

    \item \textbf{Crear un sistema de generación y almacenamiento de \gls{dataset}}: Implementar funcionalidades para la captura sincronizada de señales \acrshort{eeg} y eventos de entrada, con exportación a formatos estándar como \acrfull{hdf5}, \acrfull{edf}+ y \acrfull{bids} que faciliten el futuro entrenamiento de modelos predictivos.

    \item \textbf{Desarrollar interfaces de visualización}: Crear interfaces para monitorizar en tiempo real tanto señales \acrshort{eeg} como eventos de entrada, proporcionando múltiples modalidades de acceso (\acrfull{gui}, \acrfull{cli} y \acrfull{api} web) que permitan a usuarios con diferentes niveles de conocimiento técnico utilizar la plataforma.

    \item \textbf{Diseñar una arquitectura extensible}: Crear una arquitectura modular basada en principios \acrfull{solid}~\cite{martin2003agile} que permita la integración con herramientas existentes de \acrfull{ml} y facilite la incorporación de nuevos dispositivos y técnicas de procesamiento.
\end{enumerate}

\section{Metodología}
\label{sec:metodologia}

El desarrollo de este proyecto ha seguido una metodología incremental-iterativa que combina las ventajas de una definición temprana de requisitos con la flexibilidad necesaria para refinar el resultado en ciclos sucesivos. Esta aproximación ha permitido detectar y corregir problemas de diseño de forma temprana, mientras se mantiene la flexibilidad para adaptar la solución a medida que se profundiza en el dominio del problema.

La metodología se caracteriza por el desarrollo modular, donde cada componente del sistema se construye de forma independiente antes de proceder a su integración. Este enfoque ha facilitado la validación incremental de funcionalidades y ha permitido la identificación temprana de problemas de compatibilidad entre módulos.

Un aspecto clave del desarrollo ha sido la aplicación de técnicas de refactorización continua, especialmente la reestructuración del código siguiendo principios \acrshort{solid} para mejorar la mantenibilidad y extensibilidad del sistema. Esta decisión resultó especialmente valiosa cuando se detectaron limitaciones en los enfoques iniciales de integración entre módulos.

La validación del sistema se ha realizado mediante el desarrollo de \acrfull{api} que han permitido la verificación exhaustiva de funcionalidades y la detección de errores internos. Adicionalmente, se han integrado herramientas de calidad de código como \textit{pre-commit hooks}, tipado estático, formateo automático y verificación de documentación para mantener estándares durante todo el proceso de desarrollo.

Los detalles específicos de la metodología, incluyendo las fases del proyecto, cronograma detallado y análisis de riesgos, se presentan en el Capítulo \ref{chap:metodologia} de esta memoria.

\section{Estructura de la memoria}
\label{sec:estructura}

Esta memoria se organiza siguiendo una estructura lógica que guía al lector desde la contextualización del problema hasta la presentación de resultados y conclusiones.

El Capítulo \ref{chap:estado_arte} presenta una revisión exhaustiva del estado del arte en tecnologías \acrshort{bci}, analizando las herramientas existentes para captura y procesamiento de datos \acrshort{eeg}, identificando sus limitaciones y oportunidades de mejora que justifican el desarrollo de la solución propuesta.

El Capítulo \ref{chap:metodologia} describe en detalle la metodología de desarrollo empleada, incluyendo la planificación del proyecto, cronograma, análisis de riesgos y estrategias de mitigación. Se documentan las fases del desarrollo incremental-iterativo y las lecciones aprendidas durante el proceso.

El Capítulo \ref{chap:analisis_diseno} aborda el análisis y diseño del sistema, presentando los requisitos funcionales y no funcionales identificados, la arquitectura general de la solución y el diseño detallado de cada componente. Se justifican las decisiones de diseño principales y se describe cómo la arquitectura propuesta aborda las limitaciones identificadas.

El Capítulo \ref{chap:implementacion} documenta la implementación del sistema, describiendo las decisiones técnicas específicas, los desafíos encontrados durante el desarrollo y las soluciones adoptadas. Se presentan los aspectos más técnicos de cada módulo y se analizan las optimizaciones realizadas para garantizar el rendimiento en tiempo real.

Finalmente, el Capítulo \ref{chap:conclusiones} presenta las conclusiones del trabajo, evaluando el cumplimiento de los objetivos propuestos, analizando las limitaciones de la solución desarrollada y proponiendo líneas de trabajo futuro que podrían extender y mejorar la plataforma.

Los anexos complementan la memoria principal con documentación técnica detallada, manuales de usuario y material adicional que facilita la comprensión y utilización del sistema desarrollado.
 \chapter{Estado del arte y tecnologías}
\label{chap:estado_arte}

\lettrine{E}{n} este capítulo se presenta una revisión exhaustiva del estado actual de las tecnologías en el campo de las interfaces cerebro-computadora (BCI), analizando las herramientas existentes para captura y procesamiento de datos EEG, así como las tecnologías empleadas en el desarrollo del framework propuesto.

\section{Interfaces cerebro-computadora: panorama actual}
\label{sec:bci_panorama}

Las interfaces cerebro-computadora han experimentado un crecimiento exponencial en las últimas dos décadas, impulsadas por avances en neurociencia, procesamiento de señales y aprendizaje automático. Los sistemas BCI modernos se clasifican principalmente según el método de adquisición de señales neurológicas:

\subsection{Técnicas de adquisición de señales cerebrales}
\label{subsec:tecnicas_adquisicion}

\begin{itemize}
    \item \textbf{Electroencefalografía (EEG)}: Método no invasivo que registra la actividad eléctrica del cerebro mediante electrodos colocados en el cuero cabelludo. Ofrece excelente resolución temporal (milisegundos) pero limitada resolución espacial~\cite{eeg_bci_review}.
    
    \item \textbf{Magnetoencefalografía (MEG)}: Detecta campos magnéticos producidos por la actividad neuronal. Proporciona mejor resolución espacial que EEG pero requiere equipamiento extremadamente costoso.
    
    \item \textbf{Espectroscopía funcional de infrarrojo cercano (fNIRS)}: Mide cambios en la oxigenación sanguínea cerebral. Más robusto ante artefactos de movimiento pero con menor resolución temporal.
    
    \item \textbf{Técnicas invasivas}: Incluyen electrocorticografía (ECoG) y registros intracraneales, que ofrecen la mejor resolución espacio-temporal pero conllevan riesgos quirúrgicos.
\end{itemize}

Para el desarrollo de aplicaciones BCI no médicas, EEG se ha establecido como la técnica predominante debido a su naturaleza no invasiva, coste relativamente bajo y disponibilidad de dispositivos comerciales accesibles~\cite{eeg_bci_review}.

\section{Sistemas de adquisición EEG existentes}
\label{sec:sistemas_eeg}

\subsection{Hardware de adquisición EEG}
\label{subsec:hardware_eeg}

El mercado de dispositivos EEG se segmenta en tres categorías principales:

\subsubsection{Sistemas clínicos de alta gama}
Dispositivos como los sistemas de Neuroscan, Brain Products o ANT Neuro están diseñados para uso hospitalario y de investigación avanzada. Características:
\begin{itemize}
    \item 64-256 canales con resolución de 24 bits
    \item Frecuencias de muestreo de hasta 20 kHz
    \item Impedancia de entrada > 100 MΩ
    \item Coste: €50,000-€200,000
\end{itemize}

Ejemplos representativos incluyen Brain Products actiCAP~\cite{brainproducts_acticap} (hasta 160 canales) y g.tec g.SCARABEO~\cite{gtec_scarabeo} (hasta 256 canales), utilizados principalmente en investigación médica avanzada.

\subsubsection{Sistemas de investigación de gama media}
Incluyen dispositivos como g.tec Unicorn, OpenBCI Cyton, y Emotiv EPOC X:
\begin{itemize}
    \item 8-32 canales con resolución de 16-24 bits
    \item Frecuencias de muestreo de 250-1000 Hz
    \item Conectividad inalámbrica (Bluetooth/WiFi)
    \item Coste: €1,000-€10,000
\end{itemize}

La tabla~\ref{tab:dispositivos_gama_media} presenta una comparativa básica de los principales dispositivos en esta categoría.

\begin{table}[hp!]
\centering
\resizebox{\textwidth}{!}{
\begin{tabular}{l|c|c|c|c}
\rowcolor{udcpink!25}
\textbf{Dispositivo} & \textbf{Canales} & \textbf{Frecuencia} & \textbf{Resolución} & \textbf{Precio} \\
\hline
\textbf{g.tec Unicorn}~\cite{gtec_unicorn} & 8 & 250 Hz & 24 bits & €1,200 \\
\textbf{Emotiv EPOC X}~\cite{emotiv_epoc_x} & 14 & 128/256 Hz & 14/16 bits & €999 \\
\textbf{OpenBCI Cyton}~\cite{openbci_cyton} & 8 & 250 Hz & 24 bits & €947 \\
\textbf{Neurosity Crown}~\cite{neurosity_crown} & 8 & 256 Hz & 16 bits & €899 \\
\textbf{Neuroelectrics Enobio}~\cite{neuroelectrics_enobio} & 32 & 500 Hz & 24 bits & €5,000 \\
\end{tabular}
}
\caption{Comparativa de dispositivos EEG de investigación de gama media}
\label{tab:dispositivos_gama_media}
\end{table}

También se incluyen dispositivos avanzados como:
\begin{itemize}
    \item \textbf{Emotiv FLEX 2 Saline}~\cite{emotiv_flex_saline}: Hasta 32 canales configurables, €1,999
    \item \textbf{OpenBCI Ultracortex Mark IV}~\cite{openbci_ultracortex}: 8-16 canales modulares
    \item \textbf{Advanced Brain Monitoring B-Alert X10}~\cite{abm_b_alert}: 9 canales para aplicaciones de monitoreo
\end{itemize}

\subsubsection{Dispositivos de consumo}
Dispositivos como Muse, NeuroSky, y Emotiv Insight dirigidos a aplicaciones de biofeedback y meditación:
\begin{itemize}
    \item 4-14 canales con resolución de 12-16 bits
    \item Frecuencias de muestreo de 128-512 Hz
    \item Enfoque en facilidad de uso
    \item Coste: €200-€1,500
\end{itemize}

Dispositivos representativos incluyen:
\begin{itemize}
    \item \textbf{InteraXon Muse 2}~\cite{muse_2}: 4 canales (TP9, AF7, AF8, TP10), 256 Hz, orientado a meditación
    \item \textbf{Emotiv Insight}~\cite{emotiv_insight}: 5 canales (AF3, AF4, T7, T8, Pz), 128 Hz, €499
    \item \textbf{NeuroSky MindWave}~\cite{neurosky_mindwave}: 1 canal, básico para aplicaciones educativas, €99
    \item \textbf{FocusCalm}~\cite{focuscalm}: Dispositivo de meditación con suscripción mensual
    \item \textbf{Mendi}~\cite{mendi}: Dispositivo de entrenamiento mental con neurofeedback
\end{itemize}

\subsection{Librerías de abstracción de hardware}
\label{subsec:librerias_hardware}

\subsubsection{BrainFlow}
BrainFlow~\cite{brainflow} es una librería multiplataforma que proporciona una API unificada para más de 20 fabricantes de dispositivos EEG. Características principales:

\begin{itemize}
    \item \textbf{Abstracción de hardware}: API consistente independiente del fabricante
    \item \textbf{Soporte multiplataforma}: Windows, macOS, Linux, Android
    \item \textbf{Múltiples lenguajes}: Python, C++, C\#, Java, R, MATLAB
    \item \textbf{Streaming en tiempo real}: Acceso a datos con latencia < 50ms
    \item \textbf{Filtrado básico}: Implementaciones optimizadas de filtros digitales
\end{itemize}

Sin embargo, BrainFlow presenta limitaciones significativas:
\begin{itemize}
    \item Ausencia de interfaces gráficas de usuario
    \item Funcionalidades de procesamiento limitadas
    \item Falta de herramientas de sincronización con eventos externos
    \item Documentación limitada para casos de uso avanzados
\end{itemize}

\subsubsection{Lab Streaming Layer (LSL)}
LSL~\cite{lsl} es un protocolo y conjunto de librerías para streaming de datos de medidas temporales en experimentos de investigación:

\begin{itemize}
    \item \textbf{Protocolo de red}: Descubrimiento automático de streams en red local
    \item \textbf{Sincronización temporal}: Corrección automática de deriva de reloj
    \item \textbf{Múltiples tipos de datos}: Señales EEG, marcadores, audio, video
    \item \textbf{Robustez}: Tolerancia a interrupciones de red y reconexión automática
\end{itemize}

Limitaciones de LSL:
\begin{itemize}
    \item Requiere programación personalizada sustancial
    \item Dependencia de software propietario para muchos dispositivos
    \item Ausencia de herramientas de visualización integradas
    \item Curva de aprendizaje pronunciada para usuarios no técnicos
\end{itemize}

\section{Software de procesamiento de señales EEG}
\label{sec:software_procesamiento}

\subsection{Entornos de análisis científico}
\label{subsec:entornos_analisis}

\subsubsection{EEGLAB}
EEGLAB~\cite{eeglab} es un toolbox de MATLAB ampliamente utilizado en neurociencia:

\begin{itemize}
    \item \textbf{Procesamiento completo}: Desde preprocesamiento hasta análisis estadístico
    \item \textbf{Análisis de componentes independientes}: Implementación robusta de ICA
    \item \textbf{Visualización avanzada}: Topografías, espectrogramas, mapas 3D
    \item \textbf{Extensibilidad}: Amplio ecosistema de plugins
\end{itemize}

Limitaciones:
\begin{itemize}
    \item Dependencia de MATLAB (licencia comercial costosa)
    \item Orientado a análisis offline, no tiempo real
    \item Interfaz de usuario no intuitiva para principiantes
    \item Limitaciones en el manejo de datasets grandes
\end{itemize}

\subsubsection{MNE-Python}
MNE-Python~\cite{mne_python} es una librería de Python para análisis de señales MEG y EEG:

\begin{itemize}
    \item \textbf{Open source}: Completamente gratuito y de código abierto
    \item \textbf{API moderna}: Diseño orientado a objetos con documentación exhaustiva
    \item \textbf{Formatos estándar}: Soporte nativo para BIDS, EDF+, FIF
    \item \textbf{Análisis estadístico}: Integración con scikit-learn y scipy
    \item \textbf{Visualización}: Gráficos interactivos con matplotlib
\end{itemize}

Ecosistema MNE completo incluye:
\begin{itemize}
    \item \textbf{MNE-MATLAB}~\cite{mne_matlab}: Herramientas MATLAB para usar el software MNE
    \item \textbf{MNE-C}~\cite{mne_c}: Núcleo en C para procesamiento de señales MEG y EEG
    \item \textbf{Utilidades adicionales}: Diversos repositorios especializados para casos de uso específicos
\end{itemize}

Fortalezas de MNE-Python que lo convierten en una elección popular:
\begin{itemize}
    \item Ecosistema Python: Integración natural con NumPy, SciPy, matplotlib
    \item Comunidad activa: Desarrollo continuo y soporte comunitario
    \item Flexibilidad: Desde análisis básico hasta investigación avanzada
    \item Rendimiento: Operaciones vectorizadas optimizadas
\end{itemize}

\subsection{Herramientas de desarrollo específicas}
\label{subsec:herramientas_desarrollo}

\subsubsection{OpenViBE}
OpenViBE~\cite{openvibe} es una plataforma para diseño, prueba e implementación de aplicaciones BCI:
\begin{itemize}
    \item Entorno de desarrollo visual (drag-and-drop)
    \item Soporte para tiempo real y análisis offline
    \item Interfaz modular basada en cajas de procesamiento
    \item Limitado a investigación, no desarrollo de productos
\end{itemize}

\subsubsection{BCI2000}
BCI2000~\cite{bci2000} es un framework para investigación BCI desarrollado en C++:
\begin{itemize}
    \item Arquitectura modular para adquisición, procesamiento y aplicación
    \item Soporte para múltiples paradigmas BCI
    \item Interfaz de configuración compleja
    \item Curva de aprendizaje pronunciada
\end{itemize}

\subsection{Otras librerías especializadas}
\label{subsec:analisis_librerias}

Además de las herramientas principales mencionadas anteriormente, existen diversas librerías especializadas que abordan aspectos específicos del procesamiento EEG:

\subsubsection{Librerías de análisis complementarias}

\begin{itemize}
    \item \textbf{pyEEG}~\cite{pyeeg}: Librería Python especializada en procesamiento continuo de datos EEG basada en MNE-Python y scikit-learn. Permite generación de funciones de respuesta temporal con señales continuas.
    
    \item \textbf{NeuroKit2}~\cite{neurokit2}: Framework Python para análisis de señales, extracción de características y visualización de datos neurofisiológicos.
    
    \item \textbf{BCILAB}~\cite{bcilab}: Plataforma MATLAB para procesamiento de señales, extracción de características y aprendizaje automático orientado a BCI (descontinuada desde 2014).
    
    \item \textbf{EEG-ExPy}~\cite{eeg_expy}: Colección de experimentos en Python más que una librería per se, compatible con múltiples dispositivos (Muse, OpenBCI, Neurosity, BrainBit, g.tec).
    
    \item \textbf{Emokit}~\cite{emokit}: Librería Python de código abierto para dispositivos Emotiv EPOC, basada en ingeniería inversa (soporte limitado a modelos anteriores a 2016).
\end{itemize}

\subsubsection{Herramientas de desarrollo específicas adicionales}

\begin{itemize}
    \item \textbf{FieldTrip}~\cite{fieldtrip}: Toolbox MATLAB para análisis y preprocesamiento de datos MEG y EEG, ampliamente utilizado en investigación neurocientífica.
    
    \item \textbf{BrainStorm}~\cite{brainstorm}: Plataforma MATLAB especializada en visualización 3D, análisis de conectividad y preprocesamiento de señales MEG, EEG y fNIRS.
    
    \item \textbf{MoBILAB}~\cite{mobilab}: Herramienta MATLAB para preprocesamiento, análisis y visualización de datos EEG móviles (desarrollo discontinuado).
    
    \item \textbf{NeuroPype}~\cite{neuropype}: Plataforma comercial Python para procesamiento de señales, visualización y análisis en tiempo real con soporte extensivo para hardware comercial.
\end{itemize}

\section{Análisis de soluciones existentes}
\label{sec:analisis_soluciones}

\subsection{Limitaciones identificadas}
\label{subsec:limitaciones}

El análisis de las herramientas existentes revela varias limitaciones críticas:

\begin{enumerate}
    \item \textbf{Fragmentación del ecosistema}: Las herramientas abordan aspectos específicos pero carecen de integración:
    \begin{itemize}
        \item BrainFlow: Excelente para adquisición, limitado para procesamiento
        \item MNE-Python: Potente para análisis, no diseñado para tiempo real
        \item LSL: Bueno para sincronización, requiere programación extensiva
    \end{itemize}

    \item \textbf{Ausencia de herramientas de captura sincronizada}: Ninguna solución proporciona captura integrada de señales EEG y eventos de entrada de usuario con sincronización temporal precisa.

    \item \textbf{Barreras de entrada técnicas}: Las herramientas existentes requieren conocimientos avanzados de programación y procesamiento de señales.

    \item \textbf{Limitaciones de interfaz de usuario}: La mayoría de herramientas carecen de interfaces gráficas intuitivas para usuarios no técnicos.

    \item \textbf{Falta de estandarización}: Formatos de datos inconsistentes entre herramientas dificultan la interoperabilidad.
\end{enumerate}

\subsection{Oportunidades identificadas}
\label{subsec:oportunidades}

Las limitaciones identificadas revelan oportunidades significativas:

\begin{itemize}
    \item \textbf{Plataforma integrada}: Necesidad de una solución que combine adquisición, procesamiento y sincronización
    \item \textbf{Interfaces múltiples}: Demanda de herramientas que soporten tanto usuarios técnicos como no técnicos
    \item \textbf{Tiempo real}: Carencia de soluciones que combinen análisis avanzado con capacidades de tiempo real
    \item \textbf{Datasets de calidad}: Necesidad de herramientas que faciliten la generación de datasets bien estructurados para BCI
\end{itemize}

\section{Tecnologías empleadas en el desarrollo}
\label{sec:tecnologias_empleadas}

\subsection{Lenguajes de programación}
\label{subsec:lenguajes}

\subsubsection{Python como lenguaje principal}
Python se seleccionó como lenguaje principal por múltiples factores:

\begin{itemize}
    \item \textbf{Ecosistema científico maduro}: NumPy, SciPy, matplotlib, scikit-learn
    \item \textbf{Librerías especializadas EEG}: MNE-Python, BrainFlow con bindings nativos
    \item \textbf{Tipado estático moderno}: Soporte para type hints desde Python 3.5+
    \item \textbf{Rendimiento adecuado}: Operaciones vectorizadas con NumPy proporcionan rendimiento suficiente para tiempo real
    \item \textbf{Comunidad activa}: Amplio soporte comunitario y documentación
\end{itemize}

\subsubsection{JavaScript para frontend web}
JavaScript ES6+ se utiliza en el módulo api\_explorer:

\begin{itemize}
    \item \textbf{Arquitectura modular}: Sistema de managers para organización del código
    \item \textbf{WebSocket nativo}: Comunicación bidireccional en tiempo real
    \item \textbf{Chart.js}: Visualización de datos científicos
    \item \textbf{Bootstrap 5}: Framework CSS responsivo
\end{itemize}

\subsection{Frameworks y librerías principales}
\label{subsec:frameworks_principales}

\subsubsection{Interfaz gráfica: PyQt5}
PyQt5 se eligió para las interfaces gráficas por:

\begin{itemize}
    \item \textbf{Rendimiento superior}: Widgets nativos del sistema operativo
    \item \textbf{Arquitectura orientada a eventos}: Señales y slots para comunicación desacoplada
    \item \textbf{Capacidades avanzadas}: Gráficos 2D/3D, threading, recursos embebidos
    \item \textbf{Multiplataforma}: Soporte nativo para Windows, macOS, Linux
    \item \textbf{Estabilidad}: Framework maduro con amplia adopción en aplicaciones científicas
\end{itemize}

Comparativa con alternativas:
\begin{itemize}
    \item \textbf{Tkinter}: Más simple pero limitado en capacidades visuales
    \item \textbf{wxPython}: Menos mantenido, documentación limitada
    \item \textbf{Kivy}: Orientado a móviles, no óptimo para aplicaciones desktop
\end{itemize}

\subsubsection{Línea de comandos: Typer + Rich}
Para las interfaces CLI se utiliza la combinación Typer + Rich:

\begin{itemize}
    \item \textbf{Typer}: Framework CLI moderno basado en type hints
    \item \textbf{Rich}: Salida de consola avanzada con tablas, colores, progreso
    \item \textbf{Autocompletado}: Generación automática de completado de comandos
    \item \textbf{Validación automática}: Basada en anotaciones de tipo Python
\end{itemize}

\subsubsection{API web: FastAPI}
FastAPI se seleccionó para el módulo api\_explorer:

\begin{itemize}
    \item \textbf{Rendimiento}: Uno de los frameworks Python más rápidos
    \item \textbf{Documentación automática}: Generación automática de OpenAPI/Swagger
    \item \textbf{Validación de datos}: Integración nativa con Pydantic
    \item \textbf{WebSocket}: Soporte nativo para comunicación bidireccional
    \item \textbf{Asíncrono}: Soporte completo para asyncio
\end{itemize}

\subsection{Librerías de procesamiento científico}
\label{subsec:librerias_cientifico}

\subsubsection{NumPy y SciPy}
Base del procesamiento numérico:

\begin{itemize}
    \item \textbf{NumPy}: Arrays multidimensionales eficientes, operaciones vectorizadas
    \item \textbf{SciPy}: Algoritmos científicos optimizados (filtros, transformadas, estadística)
    \item \textbf{Rendimiento}: Implementaciones en C/Fortran con bindings Python
    \item \textbf{Interoperabilidad}: Estándar de facto para datos científicos en Python
\end{itemize}

\subsubsection{MNE-Python para análisis avanzado}
Integración opcional pero potente:

\begin{itemize}
    \item \textbf{Conversión de datos}: Transformación BrainFlow → MNE para análisis avanzado
    \item \textbf{Visualización científica}: Topografías, PSD, espectrogramas
    \item \textbf{Análisis estadístico}: Herramientas para análisis de grupos y condiciones
    \item \textbf{Formatos estándar}: Soporte nativo para BIDS, EDF+, FIF
\end{itemize}

\subsection{Tecnologías de captura de entrada}
\label{subsec:tecnologias_captura}

\subsubsection{Captura de teclado y ratón: pynput}
\begin{itemize}
    \item \textbf{Multiplataforma}: Windows, macOS, Linux
    \item \textbf{Eventos de bajo nivel}: Acceso a eventos del sistema operativo
    \item \textbf{Precisión temporal}: Timestamps de sistema con precisión de milisegundos
    \item \textbf{Flexibilidad}: Tanto captura como simulación de eventos
\end{itemize}

\subsubsection{Controladores de juego: PySDL2}
Para soporte de controladores PS4:

\begin{itemize}
    \item \textbf{SDL2 bindings}: Acceso completo a capacidades SDL2
    \item \textbf{Soporte de controladores}: Reconocimiento automático de dispositivos
    \item \textbf{Eventos thread-safe}: Procesamiento de eventos en hilos separados
    \item \textbf{Retroalimentación háptica}: Soporte para vibración (futuras extensiones)
\end{itemize}

\subsection{Tecnologías de serialización y comunicación}
\label{subsec:serializacion}

\subsubsection{Optimización JSON: orjson}
Para comunicación web de alto rendimiento:

\begin{itemize}
    \item \textbf{Rendimiento}: 2-5x más rápido que json estándar
    \item \textbf{Compatibilidad}: Drop-in replacement para json
    \item \textbf{Tipos adicionales}: Soporte nativo para datetime, UUID, dataclasses
    \item \textbf{Validación}: Manejo robusto de datos no serializables
\end{itemize}

\subsubsection{WebSocket para tiempo real}
Comunicación bidireccional de baja latencia:

\begin{itemize}
    \item \textbf{Latencia mínima}: < 15ms para eventos de datos
    \item \textbf{Worker pools}: Distribución de carga para múltiples clientes
    \item \textbf{Throttling adaptativo}: Limitación inteligente basada en frecuencia
    \item \textbf{Fallback}: Degradación grácil en caso de sobrecarga
\end{itemize}

\subsection{Formatos de exportación de datos}
\label{subsec:formatos_exportacion}

\subsubsection{BIDS (Brain Imaging Data Structure)}
Estándar para organización de datos neurocientíficos~\cite{bids}:

\begin{itemize}
    \item \textbf{Estructura jerárquica}: Organización por sujeto/sesión/tarea
    \item \textbf{Metadatos JSON}: Descripción detallada de parámetros de adquisición
    \item \textbf{Interoperabilidad}: Compatible con herramientas de análisis principales
    \item \textbf{Comunidad}: Estándar adoptado por la comunidad neurocientífica internacional
\end{itemize}

\subsubsection{EDF/EDF+ (European Data Format)}
Formato estándar para datos médicos:

\begin{itemize}
    \item \textbf{Compatibilidad clínica}: Ampliamente usado en entornos médicos
    \item \textbf{Metadatos integrados}: Información del paciente y configuración
    \item \textbf{Soporte de anotaciones}: Marcadores temporales con EDF+
    \item \textbf{Librerías maduras}: pyedflib para lectura/escritura eficiente
\end{itemize}

\subsubsection{HDF5 (Hierarchical Data Format)}
Formato de alto rendimiento para datos científicos:

\begin{itemize}
    \item \textbf{Estructura jerárquica}: Organización flexible de datasets
    \item \textbf{Compresión}: Algoritmos de compresión integrados
    \item \textbf{Acceso parcial}: Lectura de subconjuntos sin cargar todo el archivo
    \item \textbf{Metadatos extensibles}: Atributos personalizables por dataset
\end{itemize}

\subsubsection{CSV (Comma-Separated Values)}
Formato universal para máxima compatibilidad:

\begin{itemize}
    \item \textbf{Simplicidad}: Legible por humanos y máquinas
    \item \textbf{Compatibilidad universal}: Soportado por todas las herramientas
    \item \textbf{Headers dinámicos}: Adaptación automática a tipos de datos capturados
    \item \textbf{Procesamiento eficiente}: Óptimo para análisis con pandas/R
\end{itemize}

\section{Sistema de electrodos estándar 10-20}
\label{sec:sistema_electrodos}

Para comprender adecuadamente los dispositivos EEG y su aplicación en BCI, es fundamental entender el sistema internacional de colocación de electrodos 10-20~\cite{system_10_20} y sus extensiones modernas.

\subsection{Sistema 10-20 clásico}
\label{subsec:sistema_10_20}

El sistema 10-20 establece la colocación estándar de electrodos EEG basada en referencias anatómicas:

\begin{itemize}
    \item \textbf{Referencias anatómicas}:
    \begin{itemize}
        \item Nz (Nasion): Punto en la base del hueso frontal entre los ojos
        \item Iz (Inion): Punto en la protuberancia occipital externa
        \item LPA/RPA: Puntos preauriculares izquierdo y derecho
    \end{itemize}
    
    \item \textbf{Nomenclatura de electrodos}:
    \begin{itemize}
        \item Fp: Pre-frontal (toma de decisiones, control de impulsos)
        \item F: Frontal (planificación, memoria de trabajo)
        \item C: Central (control motor, procesamiento sensorial)
        \item P: Parietal (integración sensorial, percepción espacial)
        \item T: Temporal (procesamiento auditivo, reconocimiento de patrones)
        \item O: Occipital (procesamiento visual)
    \end{itemize}
    
    \item \textbf{Numeración}: Números impares (hemisferio izquierdo), pares (derecho), 'z' (línea media)

    Esto se puede observar en la figura ~\ref{fig:10_20_system}

    \begin{figure}[hp!]
\centering
\includegraphics[width=0.9\textwidth,height=15cm,keepaspectratio]{imaxes/10_20_system.png}
\caption{Sistema 10-20}
\label{fig:10_20_system}
\end{figure}
\end{itemize}

\subsection{Sistema MCN extendido}
\label{subsec:sistema_mcn}

Para sistemas de mayor resolución, el sistema MCN (Modified Combinatorial Nomenclature) introduce posiciones intermedias:

\begin{itemize}
    \item \textbf{AF}: Entre Fp y F (actividades cognitivas superiores)
    \item \textbf{FC}: Entre F y C (integración fronto-central)
    \item \textbf{FT}: Entre F y T (percepción emocional y sensorial)
    \item \textbf{CP}: Entre C y P (integración sensorial)
    \item \textbf{TP}: Entre T y P (memoria, integración multisensorial)
    \item \textbf{PO}: Entre P y O (integración visual y sensorial)
\end{itemize}

El sistema MCN también renombra electrodos del sistema 10-20:
\begin{itemize}
    \item T3 → T7, T4 → T8
    \item T5 → P7, T6 → P8
\end{itemize}

\subsection{Relevancia para aplicaciones BCI}
\label{subsec:relevancia_bci}

La comprensión del sistema de electrodos es crucial para el desarrollo de aplicaciones BCI efectivas:

\begin{itemize}
    \item \textbf{Selección de canales}: Diferentes regiones cerebrales están asociadas con distintos tipos de actividad mental
    \item \textbf{Artefactos}: Conocer la localización ayuda a identificar y filtrar artefactos (parpadeos en Fp1/Fp2, movimientos musculares en electrodos temporales)
    \item \textbf{Paradigmas BCI}: 
    \begin{itemize}
        \item Imaginación motora: Enfoque en C3, C4, Cz
        \item P300: Electrodos parietales (Pz, P3, P4)
        \item SSVEP: Electrodos occipitales (O1, O2, Oz)
    \end{itemize}
    Esto se puede observar en la figura ~\ref{fig:mcn_system}

    \begin{figure}[hp!]
\centering
\includegraphics[width=0.9\textwidth,height=15cm,keepaspectratio]{imaxes/mcn_system.png}
\caption{Sistema MCN Extendido}
\label{fig:mcn_system}
\end{figure}
\end{itemize}

\section{Arquitecturas de software adoptadas}
\label{sec:arquitecturas_software}

\subsection{Patrones arquitectónicos principales}
\label{subsec:patrones_arquitectonicos}

\subsubsection{Arquitectura dirigida por eventos}
Patrón fundamental en todos los módulos:

\begin{itemize}
    \item \textbf{Desacoplamiento}: Componentes independientes comunicados via eventos
    \item \textbf{Extensibilidad}: Fácil adición de nuevos suscriptores
    \item \textbf{Thread-safety}: Comunicación segura entre hilos
    \item \textbf{Framework-agnostic}: Sistema de eventos independiente de Qt
\end{itemize}

\subsubsection{Principios SOLID}
Aplicación sistemática de principios de diseño orientado a objetos:

\begin{itemize}
    \item \textbf{Single Responsibility}: Cada clase tiene una responsabilidad específica
    \item \textbf{Open/Closed}: Extensible sin modificar código existente
    \item \textbf{Liskov Substitution}: Interfaces intercambiables
    \item \textbf{Interface Segregation}: Interfaces específicas en lugar de monolíticas
    \item \textbf{Dependency Inversion}: Dependencias hacia abstracciones
\end{itemize}

\subsubsection{Patrón Singleton}
Para servicios centralizados:

\begin{itemize}
    \item \textbf{Thread-safe implementation}: Double-checked locking
    \item \textbf{Lazy initialization}: Inicialización bajo demanda
    \item \textbf{Global access}: Acceso consistente desde cualquier componente
    \item \textbf{Resource management}: Gestión centralizada de recursos compartidos
\end{itemize}

\subsection{Arquitectura multi-interfaz}
\label{subsec:arquitectura_multi_interfaz}

El framework implementa una arquitectura que soporta múltiples interfaces simultáneamente:

\subsubsection{Capa de servicios compartida}
\begin{itemize}
    \item \textbf{DataAccessService}: Acceso unificado a datos en tiempo real
    \item \textbf{InteractionService}: Control centralizado de funcionalidades
    \item \textbf{Estado sincronizado}: Consistency entre interfaces
\end{itemize}

\subsubsection{Adaptadores de interfaz}
\begin{itemize}
    \item \textbf{Qt Adapter}: Thread-safe bridge para operaciones GUI
    \item \textbf{API Adapter}: Exposición de servicios via REST/WebSocket
    \item \textbf{CLI Adapter}: Acceso programático para scripting
\end{itemize}

\section{Herramientas de desarrollo y calidad}
\label{sec:herramientas_desarrollo}

\subsection{Control de calidad de código}
\label{subsec:control_calidad}

\subsubsection{Pre-commit hooks}
Automatización de verificaciones de calidad:

\begin{itemize}
    \item \textbf{black}~\cite{black}: Formateo automático de código Python
    \item \textbf{mypy}~\cite{mypy}: Verificación de tipos estáticos
    \item \textbf{pydocstyle}~\cite{pydocstyle}: Verificación de docstrings
    \item \textbf{pre-commit-hooks}~\cite{pre_commit_hooks}: Verificaciones básicas (YAML, EOF, whitespace)
\end{itemize}

\subsubsection{Tipado estático}
Uso extensivo de type hints:

\begin{itemize}
    \item \textbf{Documentación automática}: Tipos como documentación
    \item \textbf{Detección temprana de errores}: mypy catch errors en desarrollo
    \item \textbf{IDE support}: Autocompletado y navegación mejorados
    \item \textbf{Refactoring safety}: Cambios seguros con verificación de tipos
\end{itemize}

\subsection{Gestión de dependencias}
\label{subsec:gestion_dependencias}

\subsubsection{Estrategia de dependencias}
El framework adopta una estrategia de dependencias escalable:

\begin{itemize}
    \item \textbf{Core mínimo}: Funcionalidades básicas con dependencias mínimas
    \item \textbf{Dependencias opcionales}: Funcionalidades avanzadas como add-ons
    \item \textbf{Fallback graceful}: Degradación elegante si dependencias opcionales no están disponibles
    \item \textbf{Version pinning}: Versiones específicas para reproducibilidad
\end{itemize}

\section{Comparativa con soluciones existentes}
\label{sec:comparativa}

La tabla~\ref{tab:comparativa_herramientas} presenta una comparativa detallada del framework desarrollado con las principales alternativas existentes.

\begin{table}[hp!]
\centering
\resizebox{\textwidth}{!}{
\begin{tabular}{l|c|c|c|c|c}
\rowcolor{udcpink!25}
\textbf{Característica} & \textbf{Framework} & \textbf{BrainFlow} & \textbf{LSL} & \textbf{MNE-Python} & \textbf{OpenViBE} \\
\hline
\textbf{Captura EEG} & \checkmark & \checkmark & \checkmark & \ding{55} & \checkmark \\
\textbf{Captura inputs usuario} & \checkmark & \ding{55} & \ding{55} & \ding{55} & \ding{55} \\
\textbf{Sincronización temporal} & \checkmark & \ding{55} & \checkmark & \ding{55} & Limited \\
\textbf{GUI integrada} & \checkmark & \ding{55} & \ding{55} & Limited & \checkmark \\
\textbf{CLI} & \checkmark & \ding{55} & \ding{55} & \ding{55} & \ding{55} \\
\textbf{Web API} & \checkmark & \ding{55} & \ding{55} & \ding{55} & \ding{55} \\
\textbf{Procesamiento tiempo real} & \checkmark & Basic & \ding{55} & \ding{55} & \checkmark \\
\textbf{Análisis avanzado} & \checkmark & \ding{55} & \ding{55} & \checkmark & Limited \\
\textbf{Formatos estándar} & \checkmark & Basic & \ding{55} & \checkmark & Limited \\
\textbf{Multiplataforma} & \checkmark & \checkmark & \checkmark & \checkmark & \checkmark \\
\textbf{Open source} & \checkmark & \checkmark & \checkmark & \checkmark & \checkmark \\
\textbf{Facilidad de uso} & \checkmark & \ding{55} & \ding{55} & \ding{55} & Limited \\
\end{tabular}
}
\caption{Comparativa de características entre herramientas BCI}
\label{tab:comparativa_herramientas}
\end{table}

\subsection{Justificación específica del hardware seleccionado}
\label{subsec:justificacion_hardware}

Para este proyecto se ha utilizado el dispositivo g.tec Unicorn Hybrid Black, cuya selección se justifica por:

\begin{enumerate}
    \item \textbf{Especificaciones técnicas adecuadas}:
    \begin{itemize}
        \item 8 canales EEG + acelerómetro/giroscopio
        \item Frecuencia de muestreo: 250 Hz (suficiente para BCI)
        \item Resolución: 24 bits (alta precisión)
        \item Conectividad Bluetooth de baja latencia
    \end{itemize}

    \item \textbf{Soporte en BrainFlow}: Integración nativa y bien documentada

    \item \textbf{Facilidad de uso}: No requiere gel conductor, configuración sencilla

    \item \textbf{Relación calidad-precio}: Equilibrio óptimo entre prestaciones y coste para investigación académica

    \item \textbf{Disponibilidad}: Accesible a través del Centro de Investigación en TIC (CITIC) de la UDC
\end{enumerate}

\section{Justificación de elecciones tecnológicas}
\label{sec:justificacion_tecnologias}

\subsection{Decisiones de diseño principales}
\label{subsec:decisiones_diseno}

\subsubsection{Python como lenguaje principal}
La elección de Python se justifica por:

\begin{enumerate}
    \item \textbf{Ecosistema maduro}: Librerías científicas establecidas y confiables
    \item \textbf{Productividad}: Desarrollo rápido con sintaxis clara
    \item \textbf{Comunidad}: Amplio soporte y documentación en neurociencia
    \item \textbf{Rendimiento suficiente}: NumPy/SciPy proporcionan rendimiento adecuado para tiempo real
    \item \textbf{Interoperabilidad}: Fácil integración con herramientas existentes
\end{enumerate}

\subsubsection{Arquitectura modular}
El diseño modular permite:

\begin{itemize}
    \item \textbf{Desarrollo independiente}: Cada módulo evoluciona separadamente
    \item \textbf{Reutilización}: Módulos utilizables en otros proyectos
    \item \textbf{Mantenimiento}: Isolación de cambios y correcciones
    \item \textbf{Testing}: Pruebas unitarias independientes por módulo
\end{itemize}

\subsubsection{Múltiples interfaces de usuario}
La estrategia multi-interfaz responde a:

\begin{itemize}
    \item \textbf{Diferentes usuarios}: Investigadores, desarrolladores, técnicos
    \item \textbf{Diferentes contextos}: Interactivo, automatizado, remoto
    \item \textbf{Integración}: Facilita adopción en workflows existentes
    \item \textbf{Accesibilidad}: Interfaces apropiadas para diferentes niveles técnicos
\end{itemize}

\section{Resumen del capítulo}
\label{sec:resumen_estado_arte}

El análisis del estado del arte revela una fragmentación significativa en el ecosistema de herramientas BCI. Mientras que existen soluciones potentes para aspectos específicos (BrainFlow para adquisición, MNE-Python para análisis, LSL para sincronización), ninguna proporciona una solución integrada que combine:

\begin{itemize}
    \item Captura sincronizada de EEG y eventos de usuario
    \item Múltiples interfaces de usuario (GUI, CLI, Web API)
    \item Procesamiento en tiempo real con análisis avanzado
    \item Exportación a formatos estándar
    \item Facilidad de uso para usuarios no técnicos
\end{itemize}

El framework desarrollado aborda estas limitaciones proporcionando una plataforma integrada que combina las fortalezas de las herramientas existentes mientras elimina las barreras de adopción identificadas. La elección de tecnologías modernas y patrones arquitectónicos sólidos garantiza que el sistema sea extensible, mantenible y preparado para futuras evoluciones en el campo BCI.

Las tecnologías seleccionadas (Python, PyQt5, FastAPI, BrainFlow, MNE-Python) representan el estado del arte en sus respectivos dominios y proporcionan una base sólida para el desarrollo de aplicaciones BCI avanzadas. La arquitectura modular y dirigida por eventos facilita la integración con herramientas existentes mientras proporciona una experiencia de usuario coherente a través de múltiples interfaces.
 \chapter{Contido demostrativo}
\label{chap:demo}

\lettrine{E}{ntre} a introdución e as conclusións, o documento conterá
tantos capítulos como sexa preciso, sempre con coidado de non rebasar
o límite de 80 páxinas fixado polo regulamento de TFGs.

Empregaremos éste de xeito demostrativo, para ilustrar o uso de
elementos habituais que poidan ser de utilidade\footnote{Por exemplo,
  isto é unha nota a pé de páxina.}.

\section{Inclusión de imaxes}

Se precisamos imaxes no noso documento, incluirémolas do xeito que se
indica na figura~\ref{fig:exemplo} (páxina~\pageref{fig:exemplo}). Se
o facemos así, \LaTeX ubicará cada imaxe no mellor lugar posible,
lugar que pode variar a medida que o documento vaia crecendo coa
inclusión de máis texto e outros elementos (máis imaxes, táboas,
etc.).

\begin{figure}[hp!]
  \centering
  \includegraphics[width=0.75\textwidth]{imaxes/udc.png}
  \caption{Pé de imaxe descritivo}
  \label{fig:exemplo}
\end{figure}

Recoméndase almacenar os ficheiros gráficos no directorio
\texttt{imaxes}.

\subsection{Inclusión de varias sub-imaxes}

Se precisamos inserir imaxes relacionadas, pode ser apropiado
incluílas como sub-figuras, do xeito que se pode apreciar na
figura~\ref{fig:exemplo-subfiguras} coas
imaxes~\ref{fig:subfigura-rotada}
e~\ref{fig:subfigura-deformada}. Como se pode ver nos exemplos desta
sección, sempre é recomendable referirse ás imaxes pola súa
referencia, xa que dese xeito non dependemos de onde queden ubicados
os elementos en cuestión.

\begin{figure}[hp!]
  \centering
  \begin{subfigure}[c]{0.3\textwidth}
    \includegraphics[angle=45,width=\textwidth]{imaxes/udc.png}
    \caption{Pé de subimaxe rotada}
    \label{fig:subfigura-rotada}
  \end{subfigure}
  \hspace{0.1\textwidth}
  \begin{subfigure}[c]{0.3\textwidth}
    \includegraphics[width=\textwidth,height=3cm]{imaxes/udc.png}
    \caption{Pé de subimaxe deformada}
    \label{fig:subfigura-deformada}
  \end{subfigure}
  \caption{Pé de imaxe xeral}
  \label{fig:exemplo-subfiguras}
\end{figure}

\clearpage

\section{Inclusión de táboas}

Se precisamos táboas no noso documento, incluirémolas do xeito que se
indica na táboa~\ref{tab:exemplo} (páxina~\pageref{tab:exemplo}). Se
o facemos así, \LaTeX ubicará cada táboa no mellor lugar posible,
lugar que pode variar a medida que o documento vaia crecendo coa
inclusión de máis texto e outros elementos (máis imaxes, táboas,
etc.).

\begin{table}[hp!]
  \centering
  \rowcolors{2}{white}{udcgray!25}
  \begin{tabular}{c|c}
  \rowcolor{udcpink!25}
  \textbf{Título de columna} & \textbf{Outro título de columna} \\\hline
  \textit{Título de fila} & Contido de celda \\
  \textit{Título de fila} & Contido de celda \\
  \textit{Título de fila} & Contido de celda \\
  \textit{Título de fila} & Contido de celda \\
  \textit{Título de fila} & Contido de celda \\
  \textit{Título de fila} & Contido de celda \\
  \end{tabular}
  \caption{Pé de táboa descritivo}
  \label{tab:exemplo}
\end{table}

Para táboas longas que ocupan varias páxinas, como é o caso da \ref{tab:longa},
recoméndase o uso do paquete \texttt{lontable}, descomentando a liña
correspondente no ficheiro raíz do proxecto (\verb+memoria_tfg.tex+).

\rowcolors{2}{white}{udcgray!25}
\begin{longtable}{l|r|c}
  \caption{Pé descritivo dunha táboa longa}
  \label{tab:longa} \\

  \rowcolor{udcpink!25}
  \textbf{Primeira columna} & \textbf{Segunda columna} & \textbf{Terceira columna} \\\hline
  \endfirsthead

  \multicolumn{3}{c}{\tablename\ \thetable{} -- {\small \textit{(vén da páxina anterior)}}} \\
  \rowcolor{udcpink!25}
  \textbf{Primeira columna} & \textbf{Segunda columna} & \textbf{Terceira columna} \\\hline
  \endhead

  \multicolumn{3}{c}{\dotfill{\small \textit{(continúa na páxina seguinte)}}\dotfill} \\
  \endfoot

  \endlastfoot

  Texto de exemplo & abcdef ghjijklmn & 123.456778 \\
  Texto de exemplo & abcdef ghjijklmn & 123.456778 \\
  Texto de exemplo & abcdef ghjijklmn & 123.456778 \\
  Texto de exemplo & abcdef ghjijklmn & 123.456778 \\
  Texto de exemplo & abcdef ghjijklmn & 123.456778 \\
  Texto de exemplo & abcdef ghjijklmn & 123.456778 \\
  Texto de exemplo & abcdef ghjijklmn & 123.456778 \\
  Texto de exemplo & abcdef ghjijklmn & 123.456778 \\
  Texto de exemplo & abcdef ghjijklmn & 123.456778 \\
  Texto de exemplo & abcdef ghjijklmn & 123.456778 \\
  Texto de exemplo & abcdef ghjijklmn & 123.456778 \\
  Texto de exemplo & abcdef ghjijklmn & 123.456778 \\
  Texto de exemplo & abcdef ghjijklmn & 123.456778 \\
  Texto de exemplo & abcdef ghjijklmn & 123.456778 \\
  Texto de exemplo & abcdef ghjijklmn & 123.456778 \\
  Texto de exemplo & abcdef ghjijklmn & 123.456778 \\
  Texto de exemplo & abcdef ghjijklmn & 123.456778 \\
  Texto de exemplo & abcdef ghjijklmn & 123.456778 \\
  Texto de exemplo & abcdef ghjijklmn & 123.456778 \\
  Texto de exemplo & abcdef ghjijklmn & 123.456778 \\
  Texto de exemplo & abcdef ghjijklmn & 123.456778 \\
  Texto de exemplo & abcdef ghjijklmn & 123.456778 \\
  Texto de exemplo & abcdef ghjijklmn & 123.456778 \\
  Texto de exemplo & abcdef ghjijklmn & 123.456778 \\
  Texto de exemplo & abcdef ghjijklmn & 123.456778 \\
  Texto de exemplo & abcdef ghjijklmn & 123.456778 \\
  Texto de exemplo & abcdef ghjijklmn & 123.456778 \\
  Texto de exemplo & abcdef ghjijklmn & 123.456778 \\
  Texto de exemplo & abcdef ghjijklmn & 123.456778 \\
  Texto de exemplo & abcdef ghjijklmn & 123.456778 \\
  Texto de exemplo & abcdef ghjijklmn & 123.456778 \\
  Texto de exemplo & abcdef ghjijklmn & 123.456778 \\
  Texto de exemplo & abcdef ghjijklmn & 123.456778 \\
  Texto de exemplo & abcdef ghjijklmn & 123.456778 \\
  Texto de exemplo & abcdef ghjijklmn & 123.456778 \\
  Texto de exemplo & abcdef ghjijklmn & 123.456778 \\
  Texto de exemplo & abcdef ghjijklmn & 123.456778 \\
  Texto de exemplo & abcdef ghjijklmn & 123.456778 \\
  Texto de exemplo & abcdef ghjijklmn & 123.456778 \\

\end{longtable}


\section{Inclusión de código fonte}

Se precisamos incluír fragmentos de código fonte, podemos facelo, por exemplo, da
seguinte maneira:

\begin{lstlisting}[language=C]
#include <stdio.h>
#define N 10

int main()
{
  int i;

  // Isto é un comentario
  puts("Ola, mundo!");

  for (i = 0; i < N; i++)
  {
    puts("LaTeX é a ferramenta de edición ideal para profesionais da informática!");
  }

  return 0;
}
\end{lstlisting}

\section{Uso da relación de acrónimos e do glosario}

Os acrónimos edítanse no ficheiro \texttt{bibliografia/acronimos.tex}
e úsanse empregando a orde \texttt{acrlong} para obter o termo
completo (deste xeito: \acrlong{erlang}), a orde \texttt{acrshort}
para obter o acrónimo (deste xeito: \acrshort{erlang}). A primeira vez
que usamos un termo con acrónimo no documento é recomendable usar orde
\texttt{acrfull} (que produce ambas versións á vez:
\acrfull{erlang}). Os acrónimos que non se usan no documento, non
aparecen na relación que se xerar na versión PDF.

Pola súa banda, os termos do glosario edítanse no ficheiro
\texttt{bibliografia/glo\-sa\-rio.tex} e úsanse empregando a orde
\texttt{gls} (deste xeito, \gls{bytecode}) ou \texttt{Gls} (deste
xeito, \Gls{bytecode}). Ao igual que os acrónimos, os termos que non
se usan no documento, non aparecen na relación que se xera na versión
PDF.

%\include{contido/...}
 \chapter{Conclusiones}
\label{chap:conclusions}
\lettrine[]{A} continuación, se detallan los aprendizajes técnicos y arquitectónicos más relevantes, la integración de los conocimientos del grado en el desarrollo del proyecto, y las oportunidades de mejora y expansión identificadas tras su implementación.

\section{Lecciones aprendidas durante implementación}
\label{sec:lecciones_aprendidas}

\subsection{Decisiones arquitectónicas importantes}
\label{subsec:decisiones_arquitectonicas}

\subsubsection{Refactorización temprana}

La refactorización temprana basada en principios SOLID fue crucial:

\textbf{Beneficios obtenidos:}
\begin{itemize}
    \item \textbf{Improved maintainability}: Código más fácil de mantener y extender
    \item \textbf{Better testing}: Componentes aislados facilitaron testing
    \item \textbf{Reduced coupling}: Menor acoplamiento entre componentes
    \item \textbf{Enhanced debugging}: Debugging más eficiente con responsabilidades claras
\end{itemize}

\subsubsection{APIs como herramienta de desarrollo}

El desarrollo de APIs fue fundamental para el proceso:

\textbf{Valor aportado:}
\begin{itemize}
    \item \textbf{Internal testing}: APIs facilitaron testing exhaustivo
    \item \textbf{Error detection}: Revelaron errores internos no detectados
    \item \textbf{Architecture validation}: Validaron la arquitectura del sistema
    \item \textbf{External integration}: Facilitaron integración con herramientas externas
\end{itemize}

\subsection{Desafíos técnicos superados}
\label{subsec:desafios_superados}

\subsubsection{Threading y Qt}

La integración de threading con Qt presentó desafíos:

\textbf{Soluciones implementadas:}
\begin{itemize}
    \item \textbf{Qt Signals/Slots}: Comunicación thread-safe entre hilos
    \item \textbf{Worker threads}: Procesamiento en hilos separados
    \item \textbf{Event-driven architecture}: Arquitectura basada en eventos para PS4
    \item \textbf{Adapter pattern}: Adaptadores para aislar threading complexity
\end{itemize}

\subsubsection{Performance optimization}

Los datos de alta frecuencia requirieron optimizaciones específicas:

\textbf{Optimizaciones implementadas:}
\begin{itemize}
    \item \textbf{Worker pools}: Distribución de carga de procesamiento
    \item \textbf{Batch processing}: Procesamiento por lotes para eficiencia
    \item \textbf{Adaptive throttling}: Throttling inteligente basado en carga
    \item \textbf{Memory optimization}: Gestión eficiente de buffers y memoria
\end{itemize}

\subsection{Relación con competencias del grado}
\label{subsec:relacion_competencias}

\subsubsection{Aplicación de conocimientos fundamentales}

El desarrollo del framework ha integrado conocimientos de múltiples asignaturas del grado:

\textbf{Fundamentos de programación y arquitectura:}
\begin{itemize}
    \item \textbf{Diseño de Software}: Aplicación de paradigmas de programación orientada a objetos
    \item \textbf{Concurrencia y Paralelismo}: Breve pincelada a patrones orientados a eventos
    \item \textbf{Algoritmos + Gestión de interfaces}: Optimización de algoritmos de procesamiento de señales
\end{itemize}

\textbf{Desarrollo software y diseño:}
\begin{itemize}
    \item \textbf{Deseño Software}: Aplicación extensiva de principios SOLID y patrones de diseño
    \item \textbf{Bases de Datos}: Diseño de estructuras de datos para almacenamiento eficiente
    \item \textbf{Interfaces Persoa Máquina}: Desarrollo de interfaces intuitivas multi-modal
\end{itemize}

\textbf{Sistemas y redes:}
\begin{itemize}
    \item \textbf{Sistemas Operativos}: Gestión de procesos, hilos y sincronización
    \item \textbf{Concorrencia e Paralelismo}: Thread-safety y procesamiento multiproceso
    \item \textbf{Redes}: Implementación de comunicación WebSocket y protocols de red
    \item \textbf{Internet e Sistemas Distribuídos}: Arquitectura REST
\end{itemize}

\subsubsection{Competencias de mención en computación}

Las asignaturas de especialización han sido particularmente relevantes:

\textbf{Inteligencia artificial y procesamiento:}
\begin{itemize}
    \item \textbf{Sistemas Inteligentes + Aprendizaxe Automática}: Base para métricas
    \item \textbf{Visión Artificial}: Aunque no directamente aplicada a señales, se dieron multiples formas de procesar imágenes que se pueden extrapolar a señales
    \item \textbf{Procesamento de Linguaxes}: Desarrollo de CLI (parsing)
\end{itemize}

\subsection{Competencias transversales desarrolladas}
\label{subsec:competencias_transversales}

\subsubsection{Investigación y aprendizaje autónomo}

\textbf{Competencias desarrolladas:}
\begin{itemize}
    \item \textbf{Investigación bibliográfica}: Análisis del estado del arte en BCI
    \item \textbf{Aprendizaje de tecnologías}: Familiarización con nuevas librerías (BrainFlow, MNE-Python)
    \item \textbf{Resolución de problemas}
    \item \textbf{Pensamiento crítico}: Evaluación y selección de alternativas tecnológicas
\end{itemize}

\section{Líneas de trabajo futuro}
\label{sec:lineas_futuro}

\subsection{Mejoras técnicas identificadas}
\label{subsec:mejoras_tecnicas}

\subsubsection{Modo headless completo}

Una limitación importante identificada es la dependencia de Qt para el funcionamiento headless:

\textbf{Problemas actuales:}
\begin{itemize}
    \item \textbf{Dependencias Qt}: CLI y APIs requieren componentes Qt inicializados
    \item \textbf{Limitaciones de deployment}: Imposibilita uso en servidores sin display
    \item \textbf{Overhead de recursos}: Inicialización innecesaria de componentes gráficos
\end{itemize}

\textbf{Soluciones propuestas:}
\begin{itemize}
    \item \textbf{Refactorización de servicios}: Separar lógica de negocio de componentes Qt
    \item \textbf{Alternative backends}: Implementación de backends alternativos sin Qt
\end{itemize}

\subsubsection{Optimizaciones de rendimiento}

\textbf{Áreas de mejora identificadas:}
\begin{itemize}
    \item \textbf{Tree.js integration}: Migración de Chart.js a Tree.js para mejor rendimiento
    \item \textbf{CUDA acceleration}: Aceleración GPU para operaciones ICA y wavelet
    \item \textbf{Compression algorithms}: Compresión en tiempo real para reducir overhead de red
\end{itemize}

 %%%%%%%%%%%%%%%%%%%%%%%%%%%%%%%%%%%%%%%%
 % Apéndices, glosarios e bibliografía  %
 %%%%%%%%%%%%%%%%%%%%%%%%%%%%%%%%%%%%%%%%

 \appendix
 \appendixpage
 \chapter{Material adicional}
\label{chap:adicional}

\lettrine{E}{xemplo} de capítulo con formato de apéndice, onde se pode
incluír material adicional que non teña cabida no corpo principal do
documento, suxeito á limitación de 80 páxinas establecida no
regulamento de TFGs.

\Blindtext

%\include{anexos/...}

 \printglossary[type=\acronymtype,title=\nomeglosarioacronimos]
 \printglossary[title=\nomeglosariotermos]

 \bibliographystyle{IEEEtranN}
 \bibliography{\bibconfig,bibliografia/bibliografia}
 \clearpage
 
\end{document}

%%%%%%%%%%%%%%%%%%%%%%%%%%%%%%%%%%%%%%%%%%%%%%%%%%%%%%%%%%%%%%%%%%%%%%%%%%%%%%%%
